% !TeX root = ../tesis.tex
%%%%%%%%%%%%%%%%%%%%%%%%%%%%%%%%%%%%%%%%%%%%%%%%%%%%%%%%%%%%%%%%%%%
% SECCIÓN 7.3 — EJE ACADÉMICO Y TRABAJO FUTURO (Issue #23)
% Papel en el hilo: lo que falta medir, y cómo unir estructura y
% demanda. Esquema: aporte Garibelt → L1 consolidar Garibelt →
% L2 otras redes (país y extranjero) → L3 calibrar lo medido →
% L4 medir lo que falta (doctorado) → cierre con 1985.
% Plan: CONTEXTO_CAP7.md §6. Decisiones 2026-11-04: sin IGISC 2026;
% sin nombrar ciudades concretas en L2 (sería especulativo).
%%%%%%%%%%%%%%%%%%%%%%%%%%%%%%%%%%%%%%%%%%%%%%%%%%%%%%%%%%%%%%%%%%%
\section{Eje Académico y Trabajo Futuro}
\label{sec:conclusiones-academico}

El tercer eje reúne lo que la tesis aporta al método y lo que deja abierto. Su
aporte principal es la Clasificación Garibelt, con el Espectro MJG como
entregable (§\ref{subsec:clasificacion-garibelt}). Una red tan heterogénea como
la de la \zmvm{}, con doce sistemas, 76 demarcaciones y una cobertura que cae
de forma abrupta al salir de la \cdmx{} (§\ref{sec:conclusiones-urbano}), no se
deja resumir en un solo número. Un promedio de indicadores esconde la dimensión que está mal; es
el sesgo de compensabilidad que señalan \citet{hossain2024} y \citet{deng2023}
(§\ref{subsubsec:garibelt-compensabilidad}). La Clasificación normaliza para
calibrar, no para sumar, como el perfil por ciudad de \citet{kujala2018}, y
asigna además a cada nodo un intervalo dominante. Así, la red actual sirvió de
línea base y los escenarios se compararon contra ella: diez nodos cambian de
intervalo dominante en el Caso~2 y 21 en el Caso~3 (§\ref{subsec:anillo-espectro-nodo}).

Se trata, con todo, de un instrumento en desarrollo, y la primera línea de
trabajo futuro es consolidarlo. Sus cortes y cotas son supuestos. Bajar la cota
de $DI$ de 2.5 a 2.0 basta para que el $DI$ de la red actual cambie de intervalo
(§\ref{subsubsec:garibelt-escala}), y los valores de $T$ y $B(v)$ dependen de
los parámetros $\alpha$ y $\beta$ del coeficiente de fricción
(§\ref{subsec:fase2-cf}). Hace falta un análisis de sensibilidad de todos
ellos. Hay también dos medidas relativas: el intervalo dominante de cada nodo se
calcula contra el máximo de cada caso, y la $FC$ se normaliza con los cuartiles
de la propia red, de modo que describe cómo se reparte la conectividad dentro
de la \zmvm{} sin decir cómo se compara con otras
(§\ref{subsubsec:garibelt-procedimiento}). Resueltos esos puntos, la
Clasificación podría servir no solo para diagnosticar una red heterogénea, sino
para ordenar varias entre sí.

Ahí empieza la segunda línea. La tesis midió una sola red real frente a dos
casos hipotéticos construidos sobre ella, así que sus resultados solo se
comparan consigo mismos. El paso siguiente es aplicar la cadena a otras redes:
las de otras zonas metropolitanas del país y las de ciudades de otros países
que tengan \termino{feeds} \gtfs{} y cartografía de base
(§\ref{subsec:datos-herramienta-replicable}). Con varios perfiles del Espectro
MJG, como los 25 que reunieron \citet{kujala2018}, se sabría si los valores de
la \zmvm{} son altos o bajos frente a otras ciudades. También se podrían
sustituir los cuartiles propios de $FC$ por una referencia común. Y la
replicabilidad que planteó §\ref{sec:conclusiones-tecnologico} quedaría
probada fuera de la ciudad donde se construyó la herramienta.

La tercera línea es calibrar lo que el modelo ya mide. La cobertura podría
calcularse con isócronas peatonales sobre la red vial de OpenStreetMap y
cruzarse con la población por manzana del Censo 2020, para pasar de superficie
cubierta a personas cubiertas (§\ref{subsec:indicador-cobertura}). La
penalización por transbordo podría contrastarse con los tiempos que devuelve la
Google Maps Directions API en horas pico y valle
(§\ref{subsec:penalizacion-transferencia}). Las velocidades de superficie
tendrían que compararse con velocidades observadas por ruta. La $FC$, por su
parte, podría ponderarse por capacidad o por flujo, como en la definición de
\citet{lee2008}, y contrastarse con el catálogo oficial de \textsc{cetram}
(§\ref{subsec:fuerza-capilar}).

La cuarta línea es medir lo que hoy queda fuera, y es la que da forma a una
propuesta de nivel doctoral. La robustez $\Delta E$ pide la eficiencia global
de \citet{latora2001efficient}, la remoción de Macro-Hubs completos y una
corrida vigente (§\ref{subsec:vulnerabilidad-estructural}); solo así se sabría
si la red anillo-radial resiste mejor una falla en el centro. A eso se suman la
demanda, la población, el ingreso y el costo de cada opción, que el modelo no
incorpora (§\ref{subsec:corredor-anillar-limitaciones}), y un Anillo Periférico
Exterior para la periferia actual. El núcleo de esa agenda es unir en un mismo
modelo los indicadores de estructura con la afluencia y las matrices
origen-destino.

Esa unión es también el cierre del argumento. En 1985, el Programa Maestro
preguntó cuánta gente usaría un corredor sobre el Periférico y, con esa sola
medida, lo descartó. Esta tesis preguntó qué aportaba ese corredor a la
estructura de la red y encontró efectos que la captación no registraba. Ninguna
de las dos preguntas basta por separado. La planeación por estructura no
sustituye a la planeación por demanda; la completa. Si lo que se mide define lo
que se construye, el trabajo que sigue es medir las dos cosas a la vez.
